\documentclass[a4paper,12pt]{article}
\usepackage[utf8]{inputenc}
\usepackage[T1]{fontenc}
\usepackage[french]{babel}
\usepackage{geometry}
\geometry{left=2.5cm,right=2.5cm,top=2.5cm,bottom=2.5cm}

\usepackage{lmodern}
\usepackage{xcolor}
\definecolor{titleblue}{RGB}{0,51,140}
\definecolor{TITLEBLUE}{RGB}{0,51,140}

\usepackage{hyperref}

\pagestyle{empty}

\begin{document}

\begin{flushleft}
    \textbf{Ismail AISSAOUI} \\
    Ingénieur d'État en Intelligence Artificielle \\
    Chlef, Algérie \\
    ismail.aissaoui.pro@gmail.com \\
    +213 660 70 77 96 \\
    \href{https://ismail-aissaoui.vercel.app}{ismail-aissaoui.vercel.app}
\end{flushleft}

\begin{flushright}
    À l'attention du \\
    \textbf{Professeur Frédérick BÉNABEN} \\
    Laboratoire DISP \\
    Mines Albi / INSA Lyon \\
    69621 Villeurbanne, France \\
    \medskip
    Le 29 Avril 2026
\end{flushright}

\bigskip

\noindent \textbf{Objet : Candidature pour un doctorat en IA pour la Gestion de Crise \& Systèmes Résilients}

\bigskip

Monsieur le Professeur,

Passionné par la modélisation des systèmes complexes et la gestion de l'incertitude via l'intelligence artificielle, je vous sollicite pour une candidature de doctorat au sein du laboratoire DISP. Vos travaux sur l'interopérabilité, la résilience et la gestion de crise par les systèmes cyber-physiques m'ont particulièrement impressionné.

Mon travail actuel à la Sonatrach porte sur un \textbf{Jumeau Numérique (Architecture Générative Mamba-SSM)} capable de simuler des scénarios de dégradation critiques ("What-If") sur des turbines à gaz. J'ai mis au point une architecture duale (\textbf{C++ JIT xLSTM} pour l'edge et substituts génératifs pour le cloud) qui permet une réponse ultra-rapide face aux anomalies, un enjeu crucial dans la gestion de crise industrielle.

Je suis convaincu que les méthodes de physique informée (\textbf{PINN}) que j'utilise pour garantir la cohérence des prédictions peuvent être étendues à la modélisation de crises multi-domaines, où la rapidité de l'inférence et la fiabilité de la simulation sont vitales.

Je joins mon CV et serais ravi d'échanger avec vous sur la manière dont mes compétences en IA générative et systèmes temps réel peuvent contribuer à vos axes de recherche.

Je vous prie d'agréer, Monsieur le Professeur, l'expression de mes salutations les plus distinguées.

\bigskip

\begin{flushright}
    \textbf{Ismail AISSAOUI}
\end{flushright}

\end{document}
